% ============================================================================
% MAIN.TEX
% File utama buku ajar Pemrograman Game
% Dapat dikompilasi sebagai satu buku lengkap
% ============================================================================

\documentclass[12pt,a4paper]{book}

% Load semua package dan konfigurasi dari preamble
% ============================================================================
% PREAMBLE.tex (REVISED & OPTIMIZED)
% Buku Ajar Pemrograman Game
% ============================================================================

% ============================================================================
% BASIC SETTINGS
% ============================================================================

% Encoding & language
\usepackage[T1]{fontenc}
\usepackage[indonesian]{babel}
% inputenc optional (pdfLaTeX legacy)
\usepackage[utf8]{inputenc} % boleh dihapus jika pakai XeLaTeX/LuaLaTeX

% Font: Times New Roman sesuai standar buku ajar Indonesia
\usepackage{times} % Times New Roman font family

% Geometry - Sesuai standar buku ajar Indonesia
% Margin: atas 2.5cm, kiri 3cm, bawah 2cm, kanan 2cm
\usepackage[a4paper, left=3cm, right=2cm, top=2.5cm, bottom=2cm]{geometry}

% Typography improvement (expansion=false: avoids pdfTeX error with non-scalable fonts)
\usepackage[expansion=false]{microtype}

% ============================================================================
% TABLE PACKAGES
% ============================================================================

\usepackage{booktabs}
\usepackage{longtable}
\usepackage{array}
\usepackage{multirow}
\usepackage{tabularx}
\usepackage{colortbl}

% ============================================================================
% MATH & ALGORITHM
% ============================================================================

\usepackage{amsmath,amssymb,amsthm,mathtools}
\usepackage{algorithm}
\usepackage{algpseudocode}

% === CODE LISTING ===
% Pilih salah satu: listings ATAU minted
\usepackage{listings}
% \usepackage{minted} % aktifkan jika pakai -shell-escape

% ============================================================================
% GRAPHICS & DIAGRAM
% ============================================================================

\usepackage{graphicx}
\usepackage{tikz}
\usepackage{pgfplots}
\pgfplotsset{compat=1.18}
\usetikzlibrary{shapes,arrows,positioning,calc,decorations.pathmorphing}
% TikZ styles for block diagrams (used in ch01, ch07, etc.)
\tikzset{
  block/.style={rectangle, draw, fill=blue!20, minimum height=1.2em, minimum width=2.5em, align=center, rounded corners=2pt},
  line/.style={draw}
}
\usepackage{float}

% ============================================================================
% REFERENCES & BIBLIOGRAPHY
% ============================================================================

\usepackage[autostyle=false,style=fallback]{csquotes}
\usepackage[style=apa,backend=biber,language=english]{biblatex}
\DeclareLanguageMapping{indonesian}{english-apa}
\DeclareLanguageMapping{english}{english-apa}
\addbibresource{references.bib}

% ============================================================================
% LISTS & STRUCTURE
% ============================================================================

\usepackage{enumitem}
\setlist[itemize]{leftmargin=*,topsep=0.5em}
\setlist[enumerate]{leftmargin=*,topsep=0.5em}

% === CUSTOM OBE ENVIRONMENTS (Learning Outcomes, Definisi, Contoh, Latihan, Catatan) ===
\newtheorem{definisi}{Definisi}[section]
\newtheorem{contoh}{Contoh}[section]
\newtheorem{latihan}{Latihan}[section]
\newtheorem{catatan}{Catatan}[section]
\newtheorem{tips}{Tips}[section]
\newenvironment{learningoutcomes}{%
  \par\noindent\textbf{Learning Outcomes:}\par\vspace{0.3em}%
  \begin{itemize}[leftmargin=*, itemsep=0.2em]%
}{%
  \end{itemize}%
}

\usepackage{subfiles}

% ============================================================================
% UTILITIES
% ============================================================================

\usepackage{xcolor}
\usepackage{fancyhdr}
\usepackage{titlesec}
\usepackage{tocloft}
\PassOptionsToPackage{nowarn}{tracklang}
\usepackage[nogroupskip,nolangwarn]{glossaries}
\usepackage{makeidx}
\usepackage{url}
\usepackage{verbatim}

% ============================================================================
% HYPERREF (load near the end)
% ============================================================================

\usepackage{hyperref}
\hypersetup{
    colorlinks=true,
    linkcolor=blue,
    urlcolor=cyan,
    citecolor=red,
    pdftitle={Buku Ajar Pemrograman Game},
    pdfauthor={Program Studi S1 Informatika},
    pdfsubject={Game Development dengan Unity},
    pdfkeywords={Unity, C\#, Game Development, OBE}
}

% Optional: smarter cross-references
\usepackage[nameinlink,noabbrev]{cleveref}

% ============================================================================
% NUMBERING CONFIGURATION
% ============================================================================

\renewcommand{\thechapter}{\arabic{chapter}}
\renewcommand{\thesection}{\thechapter.\arabic{section}}
\renewcommand{\thesubsection}{\thesection.\arabic{subsection}}

\counterwithin{figure}{chapter}
\counterwithin{table}{chapter}
\counterwithin{equation}{chapter}

% ============================================================================
% LISTINGS CONFIG (C#)
% ============================================================================

\lstset{
    language=[Sharp]C,
    basicstyle=\ttfamily\small,
    keywordstyle=\color{blue},
    commentstyle=\color{green!60!black},
    stringstyle=\color{red},
    numbers=left,
    numberstyle=\tiny\color{gray},
    backgroundcolor=\color{gray!10},
    frame=single,
    breaklines=true,
    tabsize=2
}

% ============================================================================
% HEADER & FOOTER
% ============================================================================

\setlength{\headheight}{14.5pt}
\pagestyle{fancy}
\fancyhf{}
\fancyhead[LE]{\leftmark}
\fancyhead[RO]{\rightmark}
\fancyfoot[C]{\thepage}

% ============================================================================
% CHAPTER TITLE FORMAT
% Sesuai standar buku ajar: Judul bab 16pt, subjudul 14pt
% ============================================================================

\titleformat{\chapter}[display]
{\normalfont\fontsize{16}{20}\bfseries\selectfont}
{\chaptertitlename\ \thechapter}{20pt}{\fontsize{16}{20}\bfseries\selectfont}
\titlespacing*{\chapter}{0pt}{50pt}{40pt}

% Format section (subjudul) 14pt
\titleformat{\section}
{\normalfont\fontsize{14}{18}\bfseries\selectfont}
{\thesection}{1em}{}

% Format subsection 12pt bold
\titleformat{\subsection}
{\normalfont\fontsize{12}{15}\bfseries\selectfont}
{\thesubsection}{1em}{}

% ============================================================================
% GLOSSARY & INDEX
% ============================================================================

\makeglossaries
\setacronymstyle{long-short}
\loadglsentries{glossary}

\makeindex

% ============================================================================
% SPACING
% Sesuai standar buku ajar: spasi 1.5, font 12pt
% ============================================================================

\setlength{\parindent}{1.5em}
\setlength{\parskip}{0.5em}
\renewcommand{\baselinestretch}{1.5} % Spasi 1.5 sesuai standar

\raggedbottom
\sloppy
\emergencystretch=3em


% ============================================================================
% DOKUMEN MULAI
% ============================================================================

\begin{document}

% Front matter (halaman awal)
\frontmatter

% ============================================================
% Halaman Sampul
% ============================================================
\begin{titlepage}
  \centering
  \vspace*{1cm}
  {\LARGE\bfseries Pemrograman Game\par}
  \vspace{0.5cm}
  {\large Menggunakan Unity Engine dan C\#\par}
  \vspace{0.5cm}
  {\large Buku Ajar Berbasis Outcome-Based Education (OBE)\par}
  \vspace{1cm}
  {\large Mata Kuliah: Pemrograman Game\par}
  {\large Kode: IF-XXX\par}
  \vspace{2cm}
  {\large Program Studi S1 Informatika\par}
  {\large Fakultas Teknologi Informasi\par}
  {\large Universitas Lorem Ipsum\par}
  \vfill
  {\large 2026\par}
\end{titlepage}

\cleardoublepage

% ============================================================
% Kata Pengantar
% ============================================================
\chapter*{Kata Pengantar}
\addcontentsline{toc}{chapter}{Kata Pengantar}

Puji syukur kehadirat Tuhan Yang Maha Esa atas terselesaikannya buku ajar \textit{Pemrograman Game} ini. Buku ini disusun dengan pendekatan \textbf{Outcome-Based Education (OBE)}, yang berfokus pada pencapaian kompetensi terukur sesuai dengan Capaian Pembelajaran Lulusan (CPL) dan Capaian Pembelajaran Mata Kuliah (CPMK).

Pemrograman Game merupakan bidang yang sangat menarik dan relevan dalam era digital saat ini. Unity Engine sebagai salah satu game engine terpopuler menyediakan platform yang powerful untuk mengembangkan berbagai jenis game, dari game sederhana hingga game AAA. Dengan kombinasi Unity dan C\#, mahasiswa dapat mempelajari konsep-konsep fundamental game development sekaligus mengasah keterampilan programming.

Buku ini dirancang untuk mendukung pembelajaran mahasiswa dengan pendekatan student-centered learning. Setiap bab dilengkapi dengan:
\begin{itemize}
  \item Sub-CPMK yang jelas dan terukur
  \item Materi pokok dengan contoh praktis di Unity
  \item Aktivitas pembelajaran yang mendorong eksplorasi mandiri
  \item Latihan dan refleksi untuk penguatan pemahaman
  \item Asesmen untuk mengukur pencapaian kompetensi
  \item Checklist kompetensi untuk self-assessment
\end{itemize}

Kami berharap buku ini dapat menjadi panduan yang efektif dalam perjalanan pembelajaran Anda menguasai Pemrograman Game dengan Unity Engine.

\vspace{1cm}
\begin{flushright}
Penyusun\\
Dr. Lorem Ipsum, M.Kom.
\end{flushright}

\cleardoublepage

% ============================================================
% Cara Menggunakan Buku Ini
% ============================================================
\chapter*{Cara Menggunakan Buku Ini}
\addcontentsline{toc}{chapter}{Cara Menggunakan Buku Ini}

Buku ajar ini dirancang dengan pendekatan OBE untuk memaksimalkan pencapaian pembelajaran Anda. Berikut panduan penggunaan buku ini:

\section*{Struktur Buku}

\textbf{Bab I: Pendahuluan dan Orientasi}\\
Memperkenalkan tujuan buku, keterkaitan dengan RPS, dan konteks kurikulum OBE.

\textbf{Bab II: Landasan Teori}\\
Menyajikan fondasi teoretis game development yang menjadi basis pembelajaran seluruh bab berikutnya.

\textbf{Bab III-XV: Unit Materi Inti}\\
Setiap bab mencakup satu topik utama game development dengan struktur lengkap: Sub-CPMK, materi, aktivitas, latihan, asesmen, dan checklist.

\textbf{Bab XVI: Evaluasi dan Refleksi Akhir}\\
Berisi asesmen komprehensif dan panduan refleksi untuk mengukur pencapaian kompetensi secara menyeluruh.

\textbf{Lampiran}\\
Menyediakan rubrik penilaian, glosarium istilah game development, dan referensi tambahan.

\section*{Komponen dalam Setiap Bab}

\begin{enumerate}
  \item \textbf{Sub-CPMK}: Baca dengan seksama untuk memahami kompetensi yang harus dicapai
  \item \textbf{Materi Pokok}: Pelajari dengan cermat, praktikkan semua contoh di Unity
  \item \textbf{Aktivitas Pembelajaran}: Lakukan secara mandiri atau berkelompok
  \item \textbf{Latihan}: Kerjakan untuk menguji pemahaman Anda
  \item \textbf{Asesmen}: Gunakan untuk mengukur pencapaian Sub-CPMK
  \item \textbf{Checklist}: Centang setelah yakin menguasai setiap indikator
\end{enumerate}

\section*{Tips Belajar Efektif}

\begin{itemize}
  \item Jangan hanya membaca, praktikkan semua contoh di Unity Editor
  \item Gunakan Unity Hub untuk mengelola versi dan proyek
  \item Kerjakan latihan sebelum melihat solusi
  \item Diskusikan konsep yang sulit dengan teman atau dosen
  \item Manfaatkan checklist untuk self-assessment berkala
  \item Kerjakan proyek mini untuk mengintegrasikan konsep yang dipelajari
\end{itemize}

\cleardoublepage

% ============================================================
% Identitas Mata Kuliah
% ============================================================
\chapter*{Identitas Mata Kuliah}
\addcontentsline{toc}{chapter}{Identitas Mata Kuliah}

\begin{tabular}{ll}
  Nama Program Studi & : S1 Informatika \\
  Nama Mata Kuliah & : Pemrograman Game \\
  Kode Mata Kuliah & : IF-XXX \\
  Semester & : 5 (Lima) \\
  SKS / Bobot Kredit & : 3 SKS (2 Teori, 1 Praktikum) \\
  Dosen Pengampu & : Dr. Lorem Ipsum, M.Kom. \\
  Tanggal Penyusunan & : 7 Februari 2026 \\
\end{tabular}

\vspace{1cm}

\section*{Capaian Pembelajaran Lulusan (CPL)}

CPL yang dibebankan pada mata kuliah ini mencakup kompetensi lulusan dalam aspek pengetahuan, keterampilan, dan sikap:

\begin{enumerate}
  \item \textbf{CPL-1 (Pengetahuan):} Menguasai konsep teoretis bidang pengetahuan tertentu secara umum dan konsep teoretis bagian khusus dalam bidang pengetahuan tersebut secara mendalam, serta mampu memformulasikan penyelesaian masalah prosedural.
  
  \item \textbf{CPL-2 (Keterampilan Umum):} Mampu menerapkan pemikiran logis, kritis, sistematis, dan inovatif dalam konteks pengembangan atau implementasi ilmu pengetahuan dan teknologi yang memperhatikan dan menerapkan nilai humaniora.
  
  \item \textbf{CPL-3 (Keterampilan Khusus):} Mampu merancang, mengimplementasikan, dan mengevaluasi game menggunakan Unity Engine dan C\# dengan mempertimbangkan standar kualitas software game.
  
  \item \textbf{CPL-4 (Sikap):} Menunjukkan sikap bertanggung jawab atas pekerjaan di bidang keahliannya secara mandiri dan mampu bekerja sama dalam tim.
\end{enumerate}

\section*{Capaian Pembelajaran Mata Kuliah (CPMK)}

Kemampuan atau kompetensi spesifik yang diharapkan mahasiswa kuasai setelah menyelesaikan mata kuliah:

\begin{enumerate}
  \item \textbf{CPMK-1:} Mahasiswa mampu memahami dan menjelaskan konsep dasar game development (game loop, game engine, GameObject, component).
  
  \item \textbf{CPMK-2:} Mahasiswa mampu menggunakan Unity Engine untuk membuat scene, mengelola GameObject, dan mengimplementasikan basic gameplay mechanics.
  
  \item \textbf{CPMK-3:} Mahasiswa mampu mengimplementasikan game logic menggunakan C\# dalam Unity dengan menerapkan konsep OOP dan design patterns.
  
  \item \textbf{CPMK-4:} Mahasiswa mampu mengimplementasikan sistem input, physics, animation, UI, dan audio dalam game.
\end{enumerate}

\section*{Matriks Keterkaitan CPL dan CPMK}

\begin{table}[h]
\centering
\begin{tabular}{cccc}
  \toprule
  \textbf{CPL} & \textbf{CPMK} & \textbf{Kontribusi} & \textbf{Keterangan} \\
  \midrule
  CPL-1 & CPMK-1 & Tinggi & Penguasaan konsep teoretis game development \\
  CPL-1 & CPMK-2 & Tinggi & Kemampuan menggunakan Unity Engine \\
  CPL-2 & CPMK-3 & Tinggi & Penerapan pemikiran sistematis dalam implementasi C\# \\
  CPL-2 & CPMK-4 & Sedang & Pemikiran kritis dalam implementasi sistem game \\
  CPL-3 & CPMK-2 & Tinggi & Penggunaan tools untuk pengembangan game \\
  CPL-3 & CPMK-3 & Tinggi & Implementasi dengan standar kualitas \\
  CPL-3 & CPMK-4 & Tinggi & Evaluasi kualitas implementasi game \\
  CPL-4 & CPMK-3 & Sedang & Tanggung jawab dalam implementasi kode \\
  CPL-4 & CPMK-4 & Sedang & Kerja sama dalam pengembangan game \\
  \bottomrule
\end{tabular}
\caption{Matriks Keterkaitan CPL dan CPMK}
\end{table}

\cleardoublepage

% ============================================================
% Daftar Isi
% ============================================================
\phantomsection
\addcontentsline{toc}{chapter}{Daftar Isi}
\tableofcontents
\cleardoublepage

\clearpage
\phantomsection
\addcontentsline{toc}{chapter}{Daftar Gambar}
\listoffigures
\cleardoublepage

\phantomsection
\addcontentsline{toc}{chapter}{Daftar Tabel}
\listoftables
\cleardoublepage

\phantomsection
\addcontentsline{toc}{chapter}{Daftar Kode Program}
\lstlistoflistings
\cleardoublepage


% Main matter (bab-bab utama)
\mainmatter

% Bab 1: Pendahuluan dan Orientasi Buku
\subfile{chapters/bab-01/bab-01}

% Bab 2: Landasan Teori dan Konsep Dasar
\subfile{chapters/bab-02/bab-02}

% Bab 3: Sub-CPMK 1 - Konsep Dasar Game Development
\subfile{chapters/bab-03/bab-03}

% Bab 4: Sub-CPMK 2 - Pengenalan Unity Engine
\subfile{chapters/bab-04/bab-04}

% Bab 5: Sub-CPMK 3 - Dasar Pemrograman C# untuk Unity
\subfile{chapters/bab-05/bab-05}

% Bab 6: Sub-CPMK 4 - GameObject, Component, Prefab, dan Scene
\subfile{chapters/bab-06/bab-06}

% Bab 7: Sub-CPMK 5 - Input System dan Player Controller
\subfile{chapters/bab-07/bab-07}

% Bab 8: Sub-CPMK 6 - Physics, Collision, dan Rigidbody
\subfile{chapters/bab-08/bab-08}

% Bab 9: Sub-CPMK 7 - Animation dan Animator
\subfile{chapters/bab-09/bab-09}

% Bab 10: Sub-CPMK 8 - UI/UX Game
\subfile{chapters/bab-10/bab-10}

% Bab 11: Sub-CPMK 9 - Audio dalam Game
\subfile{chapters/bab-11/bab-11}

% Bab 12: Sub-CPMK 10 - Level Design dan Game Mechanics
\subfile{chapters/bab-12/bab-12}

% Bab 13: Sub-CPMK 11 - AI Sederhana untuk Game
\subfile{chapters/bab-13/bab-13}

% Bab 14: Sub-CPMK 12 - Optimization dan Debugging
\subfile{chapters/bab-14/bab-14}

% Bab 15: Sub-CPMK 13 - Build dan Deployment
\subfile{chapters/bab-15/bab-15}

% Bab 16: Evaluasi dan Refleksi Akhir
\subfile{chapters/bab-16/bab-16}

% Back matter (bagian akhir)
\backmatter

% Lampiran
\chapter*{Lampiran}
\addcontentsline{toc}{chapter}{Lampiran}

% ============================================================
% Lampiran A: Rubrik Penilaian
% ============================================================
\section*{Lampiran A: Rubrik Penilaian Tugas Praktik}

\begin{table}[h]
\centering
\footnotesize
\begin{tabular}{|l|l|l|l|}
\hline
\textbf{Kriteria} & \textbf{Sangat Baik} & \textbf{Baik} & \textbf{Perlu Perbaikan} \\
\hline
Desain Scene & Struktur jelas & Cukup jelas & Belum terorganisir \\
\hline
GameObject Setup & Komponen tepat & Sebagian tepat & Banyak kesalahan \\
\hline
Script Implementation & Logic benar & Ada kekurangan & Banyak kesalahan \\
\hline
Code Quality & Naming rapi & Ada duplikasi & Banyak duplikasi \\
\hline
Game Mechanics & Sesuai desain & Ada bug minor & Tidak berjalan baik \\
\hline
\end{tabular}
\caption{Rubrik Penilaian Tugas Praktik}
\end{table}

% ============================================================
% Lampiran B: Referensi Unity API
% ============================================================
\section*{Lampiran B: Referensi Unity API Penting}

\subsection*{GameObject Methods}
\begin{itemize}
  \item \texttt{GameObject.Find(string name)}: Mencari GameObject berdasarkan nama
  \item \texttt{GameObject.FindWithTag(string tag)}: Mencari GameObject berdasarkan tag
  \item \texttt{Instantiate(Object original)}: Membuat instance baru dari object
  \item \texttt{Destroy(Object obj)}: Menghapus object dari scene
\end{itemize}

\subsection*{Transform Methods}
\begin{itemize}
  \item \texttt{transform.position}: Posisi object dalam world space
  \item \texttt{transform.rotation}: Rotasi object
  \item \texttt{transform.localScale}: Skala object
  \item \texttt{transform.Translate(Vector3 translation)}: Menggerakkan object
  \item \texttt{transform.Rotate(Vector3 eulerAngles)}: Merotasi object
\end{itemize}

\subsection*{Component Methods}
\begin{itemize}
  \item \texttt{GetComponent\u003cT\u003e()}: Mendapatkan component dari GameObject
  \item \texttt{AddComponent\u003cT\u003e()}: Menambahkan component ke GameObject
  \item \texttt{GetComponentInChildren\u003cT\u003e()}: Mencari component di child objects
\end{itemize}

\subsection*{Input Methods}
\begin{itemize}
  \item \texttt{Input.GetKey(KeyCode key)}: Cek apakah key sedang ditekan
  \item \texttt{Input.GetKeyDown(KeyCode key)}: Cek apakah key baru ditekan
  \item \texttt{Input.GetMouseButton(int button)}: Cek apakah mouse button ditekan
  \item \texttt{Input.GetAxis(string axisName)}: Mendapatkan nilai input axis
\end{itemize}


% Daftar Pustaka
\input{backmatter}

\end{document}

% ============================================================================
% END OF MAIN.TEX
% ============================================================================
